For each level the twin is re-simulated with a new forecast-error realization. Detection degrades gracefully. On the 14-bus system DT-GCN keeps 97.9\%, 96.7\% and 95.7\% F1 at $\sigma_f=1$, 4 and 8\%, so even at 8\% it matches the best raw-telemetry detector (SVM, 95.6\%). On the 118-bus system it drops from 98.3\% to 94.4\% F1, compared with 97.9\% to 90.9\% for DT-Threshold and 88.3\% to 79.7\% for DT-MLP. The graph model is therefore the most robust \emph{detector} when the twin is imperfect. Localization is more sensitive. On the 118-bus system, DT-GCN's bus-level F1 falls from 89.5\% ($\sigma_f=1\%$) to 64.6\% (4\%) and 54.4\% (8\%), approaching the twin-threshold level of about 50\%. On the 14-bus system the decline is milder (90.8\% to 77.1\%), and there DT-MLP degrades less (84.0\% at 8\%). Accurate localization in large networks therefore requires a well-calibrated twin, with nodal forecast errors of a few percent. This is a practical design target for DT-enabled control centers.
